\section{Settings and Save / Load}
\label{sec:settings}

\subsection{Settings dialog}

\menu{File $\rightarrow$ Settings\dots} opens a scrollable dialog of
five groups covering the application-wide preferences.

\begin{description}
    \item[\app{General}.]
        Settings file location (with an \app{Open} button to reveal
        the folder), auto-conversion of file paths between Linux and
        Windows on save-load round trips, the \app{Theme} selector
        --- \app{Dark (default)} or \app{Classic 98 (dark)}, a retro
        look with Windows-98-style 3D bevels, sharp corners, and a
        pixel font, applied live on OK --- the UI scale factor for
        HiDPI screens (restart required), and the
        \app{Verbose session log} toggle. A session log is
        \emph{always} written to a dated file under \code{logs/}
        (the newest 20 are kept); the toggle raises the log level to
        DEBUG on the next launch.
    \item[\app{Performance}.]
        Maximum number of CPU cores the fitter and the Auto-Fitter
        may use in parallel (default: cores minus two).
    \item[\app{Output}.]
        Default output directory, iteration mode (\app{Auto}
        numbers \code{iter\_NNN} folders automatically; \app{Manual}
        uses a custom label), default plot format (PNG / PDF / SVG),
        and default plot DPI.
    \item[\app{Fitting Defaults}.]
        Default optimiser, statistic, y-error model, and x-column for
        new Analysis blocks; whether the parameter-correlation
        matrix is shown by default. These seed new blocks only ---
        existing projects keep their settings.
    \item[\app{Plot Defaults}.]
        A \app{Global} tab sets matplotlib rcParams (line and marker
        sizes, font and tick sizes, axis line width, errorbar cap
        size, figure and save DPIs); per-plot-type sub-tabs
        (\app{Fit Plot}, \app{Walk Plot}, \app{Correlation},
        \app{Chi-sq Map}, \app{Tracker}, \app{Preview / TOF})
        override the global values for that plot kind --- the
        fit-plot tab includes count-adaptive marker sizing and the
        shade-under-fit / values-box options, and the isotope-shift
        plot has its own tab in the Output block's
        \app{Plot Options\dots} dialog. \app{Reset All Plot
        Defaults} returns every value to the built-in defaults.
\end{description}

Settings are persisted to \code{settings/settings.yaml} inside the
application folder (not an OS config directory), also reachable from
the \app{Open} button in the \app{General} group. Older locations
are migrated automatically on first run.

\subsection{Save and Load}

\denis{} uses a single YAML format for all save / load operations.
The file has one section per major tab plus shared top-level
registries:

\begin{lstlisting}
estimate:         {...}   # incl. last_run (restores plots/log on load)
preanalysis:      {...}   # projects: files, gates, models, plot options
analysis:         {...}   # projects (block pipelines) + isotope_shifts
scan_filters:     {...}   # per-scan exclusions, shared across tabs
calibrations:     {...}   # per-file voltage-calibration overrides
calibration_acks: [...]   # acknowledged calibration warnings
nist_asd:         {...}   # NIST ASD Browser state and ranked schemes
\end{lstlisting}

The File menu follows the standard editor model, backed by
dirty tracking (a fingerprint of the full save state):

\begin{description}
    \item[\app{Save} (\code{Ctrl+S}).] Writes to the currently
        loaded file with no dialog; falls through to Save As when no
        file is loaded yet.
    \item[\app{Save As\dots} (\code{Ctrl+Shift+S}).] Always opens a
        file dialog and writes every tab.
    \item[\app{Save Tab As\dots}.] Writes only the active tab's
        section, merging into an existing file so the other sections
        are preserved.
    \item[\app{Load All\dots} (\code{Ctrl+Shift+O}) /
          \app{Load Tab\dots} (\code{Ctrl+O}).] Populate every
        recognised section, or only the active tab's. The status bar
        reports which sections were loaded. Estimate runs and
        Analysis iterations recorded in the file are restored into
        the Results tab without re-running.
    \item[\app{New Window}.] Opens a second, fully independent
        \denis{} instance (it skips the single-instance guard and
        the splash screen).
\end{description}

A legacy top-level Estimate format (with \code{element},
\code{transition}, \code{isotopes}\dots\ at the root) is detected
automatically and routed into the \app{Estimate} tab so older
configuration files continue to work. If a loaded file references
data files that no longer exist, a locate-or-skip flow runs first
and rewrites relocated paths back into the YAML.

On exit, \denis{} prompts only when something actually changed since
the last save or load. The prompt offers \app{Save} (to the current
file, or Save As when none is loaded), \app{Don't Save}, and
\app{Cancel}; \app{Cancel} is the default so an accidental
\code{Enter} never discards work. An untouched or freshly-saved
session closes silently.
